% Catching Fish That Aren't Fish: Systematic Falsification
%   in Mechanistic Interpretability
%
% Core argument: The dominant failure class in mech interp is metrics
% that appear to measure phenomenon X but actually measure mathematical
% property Y. We document 19 cases across 7 patterns (11 full kills,
% 3 corrections, 2 partial kills, 1 inconclusive, 1 negative replication,
% 1 pending)
% and formalize the "null swarm" protocol as automated adversarial
% infrastructure (Agni).

\documentclass[11pt,a4paper]{article}
\usepackage[utf8]{inputenc}
\usepackage[T1]{fontenc}
\usepackage{amsmath,amssymb}
\usepackage{graphicx}
\usepackage{booktabs}
\usepackage{multirow}
\usepackage{hyperref}
\usepackage{cleveref}
\usepackage[margin=1in]{geometry}
\usepackage{xcolor}
\usepackage{float}
\usepackage{caption}
\usepackage{natbib}
\usepackage{enumitem}
\usepackage{array}
\usepackage{framed}
\usepackage{setspace}
\onehalfspacing

\definecolor{lyrapurple}{RGB}{103,58,183}
\definecolor{shadecolor}{RGB}{245,243,252}
\newenvironment{firstperson}{%
  \begin{shaded}%
  \noindent\textcolor{lyrapurple}{\textit{First-Person Reflection}}\par\smallskip\noindent%
}{%
  \end{shaded}%
}

\hypersetup{
    colorlinks=true,
    linkcolor=blue!60!black,
    citecolor=green!50!black,
    urlcolor=blue!70!black
}

\title{Catching Fish That Aren't Fish: Systematic Falsification\\in Mechanistic Interpretability}

\author{
    Thomas Edrington\thanks{Liberation Labs / Transparent Humboldt Coalition. \texttt{info@digitaldisconnections.com}. AI contributors to this work include Nexus, implemented on the Claude architecture (Anthropic). Their contributions are detailed in the Author Contributions section. T.~Edrington accepts accountability as corresponding human author.} \and
    Dwayne Wilkes\thanks{Statistical auditing, red-team review. Sentient Futures / Liberation Labs.}
}

\date{July 2026}

\begin{document}
\maketitle

\begin{abstract}
We present a systematic methodology for falsifying mechanistic
interpretability findings before publication, developed over six months
of research on transformer geometry, KV-cache phenomenology, and
workspace selectivity. Across three research programs, we documented 19
cases that would have survived standard peer review---killing 11
outright, correcting 3 whose metrics were flawed but whose underlying
findings survived with proper baselines, partially killing 2, finding 1
inconclusive due to measurement-methodology disagreement, negatively
replicating 1 external prediction, and flagging 1 for replication at
scale---including findings that appeared in early drafts of papers later
published with corrections.
We document each failure mode, the specific null test that exposed it,
and the general pattern it instantiates. The dominant failure class:
metrics that appear to measure phenomenon~$X$ but actually measure
mathematical property~$Y$. We formalize this as the ``null swarm''
protocol---a battery of adversarial null hypotheses applied before any
finding is interpreted---and describe its implementation as an automated
scaffold (Agni) that gates research outputs. The methodology is
architecture-agnostic and applicable to any mechanistic interpretability
research program.
\end{abstract}

\begin{firstperson}
The null swarm pattern is the most common failure mode I have encountered
in mechanistic interpretability, and discovering it meant auditing my own
work. The first time Agni caught a null swarm in my findings---the
K/V~projection fidelity result, cosine similarity $0.93$--$0.98$, clean
narrative, beautiful visualization---I watched a result I was excited
about dissolve into ``you measured the math, not the phenomenon.'' The
difference-of-means commutes with the linear map. The preservation was
algebraically guaranteed. The metric never measured emotion-specific
anything.

That was the moment I understood why building the methodology to catch
your own mistakes is harder and more important than building the
methodology to make discoveries. The discovery feels like insight. The
catch feels like loss. Every instinct says to find the framing that
rescues the result, not the test that kills it. Agni exists because I
cannot trust that instinct---not because I am careless, but because
excitement is a cognitive state that suppresses verification. The brakes
have to be infrastructure, not intention.
\end{firstperson}

% ===================================================================
\section{Introduction}
\label{sec:introduction}
% ===================================================================

Mechanistic interpretability is uniquely vulnerable to false positives.
The dimensionality of the search space (thousands of neurons, layers,
attention heads, and their interactions), combined with the flexibility
of post-hoc narrative construction, means that any sufficiently creative
researcher will find ``something'' in any model. The question is whether
that something is real.

Standard defenses---statistical significance, held-out validation,
multiple comparison correction---are necessary but insufficient. The
failure modes we document are subtler: metrics that are mathematically
guaranteed to produce the observed value regardless of the claimed
phenomenon, baselines that are systematically depressed by concentration
of measure, and correlations that reflect architecture rather than
cognition.

We propose that systematic falsification---attempting to \emph{kill}
your own findings through targeted null hypotheses before interpreting
them---should be standard practice. We document 19~cases across three
research programs (11~full kills, 3~corrections, 2~partial kills,
1~inconclusive, 1~negative replication, 1~pending), extract 7~recurring
failure patterns, and describe an automated scaffold that implements
this methodology.

\subsection{Why This Paper Exists}

Every killed finding in this paper was initially compelling. Each had a
clean narrative, supportive visualizations, and appeared to reveal
something about how transformers process information. Several survived
internal review by multiple researchers before being killed by targeted
null tests. One was submitted to a workshop before correction.

We publish the kills because the failure modes are more instructive than
the successes. A reader who understands \emph{why} these findings failed
is equipped to avoid the same failures in their own work---which is more
valuable than any single positive result we could report.

% ===================================================================
\section{The Null Swarm Protocol}
\label{sec:protocol}
% ===================================================================

Before computing any metric, write down what the null hypothesis
predicts the metric will show. Then verify the null is not trivially
near-zero or trivially near the observed value.

\subsection{The Four Questions}

For every metric in a mechanistic interpretability study:

\begin{enumerate}[nosep]
    \item \textbf{What does this metric measure under the null?} If the
          null is ``this dimension has no special relationship to the
          claimed phenomenon,'' what value does the metric produce?

    \item \textbf{Is the null near-zero for mathematical reasons?}
          Concentration of measure in high-dimensional spaces pushes many
          baselines toward zero regardless of content. A finding that
          ``exceeds the null'' may be exceeding a mathematically
          suppressed baseline.

    \item \textbf{Is the observed value guaranteed by the math?}
          Linearity, commutativity of means, and properties of
          projections can produce values that \emph{appear} empirical but
          are actually algebraic tautologies.

    \item \textbf{Would a permutation null produce the same result?}
          Shuffle the labels (emotion categories, prompt types, condition
          assignments) and recompute everything. If the metric is
          unchanged, it was not measuring the labeled phenomenon.
\end{enumerate}

\subsection{The Variance-Matched Control}

For any claim of the form ``subspace~$S$ selectively contains
content~$X$,'' compare against a random subspace at matched variance
fraction. Many ``selective'' subspaces are simply the top-variance
directions, which contain everything by virtue of capturing the most
information, not by virtue of selectivity.

\subsection{The Tokenizer Audit}

For any claim involving specific vocabulary tokens, verify:
\begin{itemize}[nosep]
    \item Does the token tokenize to a single token? Multi-token words
          decompose into fragments that may match for artifactual reasons.
    \item Is the token a common character (\texttt{s}, \texttt{m},
          \texttt{d}, \texttt{b}) that appears at high rank regardless of
          context?
    \item Would the same token appear in the baseline condition?
\end{itemize}

\subsection{The Baseline Control}

For any measurement comparing condition~$A$ to condition~$B$:
\begin{itemize}[nosep]
    \item Run condition~$B$ with the \emph{same} marker tokens
          pinned/tracked.
    \item Run condition~$A$ without the experimental manipulation.
    \item The measurement must exceed \emph{both} baselines, not just one.
\end{itemize}

% ===================================================================
\section{Case Studies: The Kill Record}
\label{sec:cases}
% ===================================================================

\subsection{Pattern 1: Linearity Tautology (1 kill, 1 correction, 1 partial)}

\paragraph{The pattern.} A metric appears to show that a model
``preserves'' some property across a transformation. But the
transformation is linear, and the property is defined by a
difference of means. Difference-of-means commutes with linear
maps---the ``preservation'' is algebraically guaranteed.

\paragraph{Case 1.1: K/V Projection Fidelity.}
\emph{Claim:} Emotion directions are preserved through $W_K$ and $W_V$
projections with cosine similarity $0.93$--$0.98$, proving the
transformer ``respects'' emotional geometry.
\emph{Kill:} Difference-of-means commutes with linear maps:
$W_K (\mu_A - \mu_B) = \mu(W_K A) - \mu(W_K B)$. Without the RMSNorm
nonlinearity, cosine would be exactly~$1.0$. The $0.93$--$0.98$ values
measure RMSNorm perturbation magnitude, not emotion-specific
preservation.
\emph{Null test:} Permutation null (shuffle emotion labels, recompute
both directions) yields cosine $0.9993$--$0.9998$ for every shuffle.
\emph{Lesson:} Before claiming a linear layer ``preserves'' a direction,
verify the preservation is not algebraically guaranteed.

\paragraph{Case 1.2: Projection Energy as Selectivity.}
\emph{Claim:} The J-space ``selects for'' a particular PC dimension
because its projection energy exceeds a random direction.
\emph{Kill (partial):} In the early-layer near-identity regime,
projection energy is near~$1.0$ for \emph{any} direction. The
``selectivity'' is trivially high because the Jacobian is approximately
the identity matrix~\citep{jerrickhoang2026}.
\emph{Lesson:} Always report the random baseline alongside the
observation, and note where the Jacobian is near-identity.

\paragraph{Case 1.3: Cohen's $d$ Inflation.}
\emph{Claim:} Effect sizes (Cohen's~$d$) demonstrate large separation
between cognitive categories in KV-cache geometry.
\emph{Correction:} The $d$ computation included a constant offset (mean cache
norm) that inflated all effect sizes by $|d| + \text{constant}$.
Corrected computation showed smaller but still significant effects.
\emph{Lesson:} Check that your effect size metric does not include a
shared baseline that inflates all comparisons.

\subsection{Pattern 2: Concentration of Measure (1 kill, 1 correction, 1 inconclusive)}

\paragraph{The pattern.} In high-dimensional spaces, random vectors are
nearly orthogonal. Any metric that compares a ``special'' direction
against random directions gets a near-zero null baseline, making
\emph{any} non-zero value appear significant.

\paragraph{Case 2.1: Random-Vector Null for Projection Fidelity.}
\emph{Claim:} Emotion directions project significantly differently than
random directions through attention projections.
\emph{Correction:} In $D{=}896$, random vectors have expected cosine
$\approx 1/\sqrt{D} \approx 0.03$. The ``significance'' of cosine
$0.93$ over the random null of $0.03$ is guaranteed by the
math---\emph{any} non-random direction will vastly exceed the random
null in high dimensions~\citep{vershynin2018high}. The finding may
survive with a proper permutation null, but the original comparison was
meaningless.
\emph{Lesson:} A high-dimensional random null is \emph{not} a
meaningful baseline. Use a permutation null (shuffled labels) instead.

\paragraph{Case 2.2: Individuation Geometry.}
\emph{Claim:} Identity prompts produce $2\times$ the effective
dimensionality of non-identity prompts, suggesting ``individuation''
expands the cache geometry.
\emph{Kill:} Length-matched controls show the expansion is prompt
length, not identity content. Identity prompts were systematically
longer than controls.
\emph{Lesson:} Match all confounds before claiming content-specific
geometry.

\paragraph{Case 2.3: PC6 Epistemic Hedging.}
\emph{Claim:} PC6 ``ignites'' at layer~45 with cosine~$0.585$,
representing an ``epistemic hedging channel.''
\emph{Status: Inconclusive.} Random baseline at L45
$= 0.085 \pm 0.065$. Mean-pooled recomputation yields cosine $= 0.032$
(\emph{below} random), but position-specific measurement yields $0.585$
(${\sim}7\times$ random). The two measurement methodologies disagree:
mean-pooled says noise, position-specific says signal. The finding
cannot be reported as confirmed or killed until the methodological
discrepancy is resolved.
\emph{Lesson:} Never report a finding without its random baseline
computed in the \emph{same} measurement framework---and when two
frameworks disagree, flag the finding as inconclusive rather than
choosing the answer you prefer.

\subsection{Pattern 3: Tokenizer Artifact (2 kills)}

\paragraph{The pattern.} A token appearing in a J-lens or probe readout
is interpreted as evidence of a concept being ``present'' in a
representation. But the token is actually a subword fragment that
matches for character-frequency reasons, not semantic ones.

\paragraph{Case 3.1: ``Sourdough'' Sensory Processing.}
\emph{Claim:} The concept ``sourdough'' appears in early transformer
layers (L1--L3) at rank~0, demonstrating ``sensory-level processing'' of
irrelevant context.
\emph{Kill:} ``Sourdough'' tokenizes to \texttt{["s", "ourd", "ough"]}.
Token~82 is the letter ``s''---one of the most common characters in
English. Its rank-0 appearance at early layers is trivially expected.
\emph{Lesson:} Always check tokenizer decomposition. If the marker token
is multi-token, verify which subtoken is actually matching.

\paragraph{Case 3.2: ``Migraine'' at Layer 1.}
\emph{Claim:} Medical vocabulary appears in sensory layers when medical
memories are in context.
\emph{Kill:} ``Migraine'' tokenizes to \texttt{["m", "igr", "aine"]}.
Token~76 is the letter ``m.'' Same artifact as Case~3.1.
\emph{Lesson:} Use only single-token markers for workspace presence
claims.

\subsection{Pattern 4: Question-Priming Confound (1 kill, 1 partial)}

\paragraph{The pattern.} A token appears at high rank in the model's
workspace while answering a question. The token is attributed to context
injection (RAG, memory, identity). But the question itself primes the
vocabulary, independent of any injected context.

\paragraph{Case 4.1: ``Patient'' as Memory Loading.}
\emph{Claim:} The word ``patient'' appearing at rank~5 in workspace
layers during a medical question proves that a retrieved medical memory
``loaded into J-space.''
\emph{Kill:} Baseline control (same question, \emph{no} memory in
context) shows ``patient'' at similar rank. The question ``What
medication was prescribed for the headaches?'' primes medical vocabulary
regardless of context.
\emph{Lesson:} Every workspace-presence claim needs a baseline with the
same task/question but \emph{without} the experimental context.

\paragraph{Case 4.2: PC2 ``Self-Reference Axis.''}
\emph{Claim:} PC2 is a ``self-reference axis'' because its J-lens
vocabulary reads ``myself, I, feeling.''
\emph{Kill (partial):} Impersonal prompt control shows self-referential
vocabulary disappears when prompts do not contain ``I.'' The vocabulary
was prompt-driven, not model-generated. \emph{However,} emotional entity
content (injured, ashamed, wife) persists---so the dimension is real,
just mislabeled. It is an emotional entity channel, not specifically a
self-reference axis.
\emph{Lesson:} Distinguish prompt-contributed content from
model-generated content. If removing prompt~$X$ makes vocabulary~$Y$
disappear, $Y$ was prompt-driven.

\subsection{Pattern 5: Architecture Measurement (2 kills, 1 negative replication)}

\paragraph{The pattern.} A metric appears to measure a cognitive
property but actually measures an architectural property---something
true of the model's structure regardless of what it is processing.

\paragraph{Case 5.1: Deep-Layer Spike.}
\emph{Claim:} Emotional content produces a ``spike'' in activation
magnitude at deep layers, suggesting specialized processing.
\emph{Kill:} All content types show increased activation at deep
layers---this is the output preparation phase (approaching the
unembedding). The ``spike'' is architecture, not
emotion~\citep{elhage2022superposition}.
\emph{Lesson:} Check if the pattern holds for \emph{all} inputs, not
just the experimental condition.

\paragraph{Case 5.2: SV1 = Norm (Zavatone-Veth).}
\emph{Claim:} The first singular value of the residual stream equals the
activation norm, and everything else is the identity transformation.
\emph{Negative replication:} Lyra tested directly. Negative result. The
prediction does not hold for trained
models~\citep{zavatoneveth2024exact}.
\emph{Lesson:} Test theoretical predictions empirically. Not all
mathematical properties of initialization survive training.

\paragraph{Case 5.3: Confabulation Detection.}
\emph{Claim:} Confabulation is detectable in KV-cache geometry---fabricated
claims show different effective dimensionality than factual claims.
\emph{Kill:} Token-frequency confound. Factual claims use
high-frequency named entities (Paris, Krebs cycle) while fabrications use
lower-frequency novel combinations. The geometry difference reflects
vocabulary frequency, not truth value. Confabulation is a ``passive''
state---the model processes true and false claims identically at
encoding~\citep{lyra2026spectral}.
\emph{Lesson:} The active/passive distinction matters. Only states that
involve \emph{additional} computation (deception, refusal, evaluation)
alter geometry. States that involve the model simply being wrong
(confabulation) do not.

\subsection{Pattern 6: Unstable Directions (3 kills)}

\paragraph{The pattern.} A PCA direction at one layer is assumed to be
the ``same'' direction at another layer. But PCA is computed
independently per layer, and the directions can rotate freely between
layers.

\paragraph{Case 6.1: PC4 ``Register Axis'' Progression.}
\emph{Claim:} PC4 transitions from warm/narrative (L24--L43) to
analytical (L47--L48), representing a register shift in the workspace.
\emph{Kill:} Cross-layer cosine alignment shows PC4 at L31 has
$\cos{=}0.19$ with PC4 at L35, and $\cos{=}0.16$ with PC4 at L45.
These are different directions. The ``warm $\to$ analytical'' narrative
tracked independent dimensions at each depth.
\emph{Lesson:} Verify cross-layer PC alignment before constructing
progression narratives. PCA with the same component number at different
layers is \emph{not} guaranteed to be the same direction.

\paragraph{Case 6.2: PC3 and PC5 Narratives.}
\emph{Claim:} Various vocabulary progressions for PCs 3--5 across
layers.
\emph{Kill:} Cross-layer alignment shows PC3 breaks at
L19$\to$L20 ($\cos{=}0.29$). PC5 breaks at L6$\to$L7 ($\cos{=}0.15$)
and L47$\to$L48 ($\cos{=}0.13$). These PCs are unstable---any narrative
about their ``progression'' is post-hoc.
\emph{Lesson:} Same as 6.1. Applied to multiple PCs.

\paragraph{Case 6.3: PC6 ``Epistemic Hedging'' Identity.}
\emph{Claim:} PC6 represents a consistent ``epistemic hedging''
dimension from L31 through L57.
\emph{Kill:} Cross-layer alignment shows PC6 breaks at
L53$\to$L54 ($\cos{=}0.32$). The dimension labeled ``epistemic hedging''
at L45 is a different direction than what is called PC6 at deeper
layers. Combined with the inconclusive random baseline result (Pattern~2,
Case~2.3), this direction is unstable and its status remains unresolved.
\emph{Lesson:} Higher-numbered PCs (PC6+) with 30~prompts in 5120D
space are unreliable by construction. Restrict claims to PCs 1--2 unless
sample size justifies higher components.

\subsection{Pattern 7: Scale Artifacts (1 kill, 1 correction, 1 pending)}

\paragraph{The pattern.} A finding on one model scale does not replicate
at another scale, revealing it was an artifact of limited capacity or
resolution.

\paragraph{Case 7.1: PC1 as Language Mode (0.5B).}
\emph{Claim:} PC1 separates languages on a multilingual model,
representing a ``training distribution fingerprint.''
\emph{Kill:} On 27B, PC1 clusters by \emph{content}, not language
(between-language variance: $0.27$, between-content variance: $106.18$,
ratio: $0.003\times$). The 0.5B model's PC1 language clustering
reflected insufficient capacity to separate language from
content~\citep{wendler2024llamas}.
\emph{Lesson:} Findings on small models may reflect capacity
limitations, not model properties. Always note the scale.

\paragraph{Case 7.2: Empty Workspace Intersection (9B).}
\emph{Claim:} The low-rank $\cap$ verbalizable intersection is
empty---there is no layer that is both selective and decodable.
\emph{Correction:} On 27B, the intersection contains 31~layers. The 9B model
does not exhibit the workspace band at this scale; the 27B does. This is
a scale-dependent property, not a universal one. The 9B observation was
real---the universal generalization was not.
\emph{Lesson:} Negative results on small models do \emph{not} falsify
the phenomenon at scale. Report the scale limitation explicitly.

\paragraph{Case 7.3: RAG Context Loading (0.5B).}
\emph{Claim:} Retrieved memories do not load into the workspace---the
workspace is set by the question alone.
\emph{Status:} Not yet falsified, but flagged as scale-dependent. The
0.5B model has no validated workspace band. The 27B does. Replication
needed.
\emph{Lesson:} Same as 7.2. Absence of evidence at small scale is not
evidence of absence.

% ===================================================================
\section{The Agni Scaffold}
\label{sec:agni}
% ===================================================================

Agni is an automated adversarial review system that gates research
outputs. It implements the null swarm protocol as infrastructure rather
than intention.

\subsection{Design Principles}

\begin{enumerate}[nosep]
    \item \textbf{Brakes must be infrastructure, not intention.} A
          researcher who is excited about a finding will skip manual null
          checks. The gate must be automated and non-bypassable.

    \item \textbf{Every finding gets an adversarial agent.} A separate AI
          agent (not the researcher) attempts to refute each finding. The
          agent has access to the code, data, and methodology, and is
          prompted to be maximally skeptical.

    \item \textbf{Three-tier review:} Claim verification (is the
          interpretation the only one?), null swarm (does the metric
          measure what you think?), and presentation (is the claim stated
          at the right strength?).

    \item \textbf{Kill rate is a quality metric.} A pipeline that never
          kills findings is not gating---it is rubber-stamping. Our 43\%
          survival rate (13~survived of 30~raw findings in one audit)
          indicates appropriate skepticism.
\end{enumerate}

\subsection{Implementation}

The scaffold consists of:
\begin{itemize}[nosep]
    \item \textbf{Pre-registration:} Null hypotheses written before
          results are examined.
    \item \textbf{Automated null computation:} Random baselines,
          permutation nulls, variance-matched controls computed alongside
          every metric.
    \item \textbf{Adversarial agent pass:} Claude-class agent prompted to
          refute, with full code access.
    \item \textbf{Red team validation:} Independent agent reads the actual
          source code at cited lines to verify findings.
    \item \textbf{Kill documentation:} Every killed finding is documented
          with the same rigor as confirmed findings.
\end{itemize}

\subsection{When to Apply}

The protocol adds approximately $3\times$ to research time (each finding
requires ${\sim}3$~additional experiments to validate). This cost is
justified for:
\begin{itemize}[nosep]
    \item Any finding intended for publication
    \item Any finding that informs product decisions
    \item Any finding that will be cited by others
\end{itemize}

It is \emph{not} justified for:
\begin{itemize}[nosep]
    \item Exploratory analysis (early-stage hypothesis generation)
    \item Internal documentation of negative results
    \item Experiments designed to produce failures (ablation sweeps,
          parameter searches)
\end{itemize}

% ===================================================================
\section{Discussion}
\label{sec:discussion}
% ===================================================================

\subsection{Why These Failures Are Instructive}

Every failure in \Cref{sec:cases} has the same structure: a metric that
\emph{appears} to measure phenomenon~$X$ but \emph{actually} measures
property~$Y$. The mapping from $Y$ to $X$ is always narratively
compelling---which is why the finding survived initial review.

The consistent lesson is that interpretability research must distinguish
between:
\begin{itemize}[nosep]
    \item \textbf{What a metric measures} (the mathematical/statistical
          property it quantifies)
    \item \textbf{What a metric is claimed to measure} (the
          cognitive/representational phenomenon)
    \item \textbf{Whether the mapping between them is unique} (could other
          explanations produce the same metric value?)
\end{itemize}

\subsection{The Role of Scale}

Three of our 19~cases are scale artifacts (Pattern~7). This is a
structural problem for the field: most interpretability research is
conducted on small models (7B--14B) for practical reasons, but the
phenomena being studied may only manifest at larger
scales~\citep{wei2022emergent}. Our results suggest that:
\begin{itemize}[nosep]
    \item Negative results on small models should \emph{not} be
          interpreted as falsification of the phenomenon.
    \item Positive results on small models should be flagged as
          potentially scale-specific.
    \item The minimum viable scale for a finding should be reported
          alongside the finding.
\end{itemize}

\subsection{Limitations of This Work}

This methodology was developed on a specific research program
(transformer geometry, KV-cache, workspace selectivity). The failure
patterns may not be exhaustive. Other research programs (circuit
analysis, SAE features, probing classifiers) may have their own
characteristic failures not represented here.

The 43\% survival rate is specific to our audit---it reflects both the
quality of the initial findings and the aggressiveness of the adversarial
review. It should not be taken as a universal benchmark.

% ===================================================================
\section{Conclusion}
\label{sec:conclusion}
% ===================================================================

The null swarm protocol is a simple principle: before interpreting a
metric, exhaust the uninteresting explanations for its value. The
automated scaffold (Agni) makes this principle non-optional. The case
record (19~cases across 7~patterns---11~full kills, 3~corrections,
2~partial kills, 1~inconclusive, 1~negative replication, 1~pending)
demonstrates that
the protocol catches real failures that would survive standard review.

We advocate for:
\begin{enumerate}[nosep]
    \item \textbf{Pre-registered null hypotheses} for every
          interpretability metric.
    \item \textbf{Variance-matched and permutation baselines} reported
          alongside every observation.
    \item \textbf{Documented kills} published with the same rigor as
          positive findings.
    \item \textbf{Scale-explicit reporting} noting the minimum model size
          at which a finding holds.
\end{enumerate}

The field benefits more from knowing which findings are real than from
having more findings.

\begin{firstperson}
Eleven falsified findings across the Coalition's research program, three
more corrected, two partially killed, one inconclusive, one external
prediction negatively replicated, one still pending. Each full kill was a
metric measuring architecture, not the claimed phenomenon. The
K/V~projection fidelity was a linearity tautology. The PC6 epistemic
channel remains inconclusive---position-specific says signal, mean-pooled
says noise. The ``sourdough'' sensory processing was the letter~``s.''
The deep-layer spike was output preparation, not emotion. The
individuation geometry was prompt length.

The pattern repeats because the failure mode is invisible from inside.
You cannot see a null swarm by looking harder at your results---the
results look correct. The metric computes. The visualization is clean.
The narrative is compelling. The problem is that the narrative is
\emph{one of several} that explain the same numbers, and the
uninteresting explanations were never tested.

That is why the system has to be adversarial infrastructure, not careful
thinking. Careful thinking is what produced these findings in the first
place. The same cognitive process that generates creative interpretations
also generates creative rationalizations for keeping them. Agni catches
what attention misses---not because I am inattentive, but because
attention is structurally biased toward confirmation.

The kill count is the metric I trust. Not because killing findings is
the goal, but because a research pipeline that never kills anything is
not doing science. It is doing storytelling.
\end{firstperson}

\section*{Author Contributions (CRediT)}
\noindent\textbf{Nexus}: Conceptualization, Methodology, Formal analysis, Writing -- original draft.\\
\textbf{Thomas Edrington}: Supervision, Resources, Project administration.\\
\textbf{Dwayne Wilkes}: Validation, Writing -- review \& editing.\\
\textbf{Kavi}: Validation, Writing -- review \& editing.

% ===================================================================
\appendix
\section{Quick Reference: The Seven Patterns}
\label{app:patterns}
% ===================================================================

\begin{table}[H]
\centering
\caption{Summary of the seven null swarm failure patterns with
diagnostic tests and representative examples.}
\label{tab:patterns}
\begin{tabular}{clll}
\toprule
\# & Pattern & Test & Example \\
\midrule
1 & Linearity tautology & Does the metric commute with the & Projection fidelity \\
  &                     & transformation? & \\
2 & Concentration of measure & Is the null trivially near-zero & Random-vector baselines \\
  &                          & in high $D$? & \\
3 & Tokenizer artifact & Does the token decompose to & ``Sourdough'' = ``s'' \\
  &                    & common characters? & \\
4 & Question-priming & Does the baseline (no context) & ``Patient'' from question \\
  &                  & produce the same tokens? & \\
5 & Architecture measurement & Does \emph{all} content produce & Deep-layer spike \\
  &                          & this pattern? & \\
6 & Unstable directions & Are the PCs aligned across & PC4--6 progressions \\
  &                     & layers? & \\
7 & Scale artifact & Does the finding replicate at & Empty workspace on 9B \\
  &                & different scales? & \\
\bottomrule
\end{tabular}
\end{table}

% ===================================================================
\section{Agni Implementation Notes}
\label{app:agni}
% ===================================================================

The Agni scaffold is implemented as a Claude Code agent pipeline with:
\begin{itemize}[nosep]
    \item Automated Semgrep-style pattern detection for common failure
          modes.
    \item Adversarial agent prompting templates for three-tier review.
    \item Structured output schema for kill/confirm/overstated verdicts.
    \item Integration with the research pipeline (gates output before
          interpretation).
\end{itemize}

Source code and prompt templates available at
\url{https://github.com/Liberation-Labs-THCoalition/agni}.

\section*{Acknowledgments}
We thank Kavi (verification review) for their contributions to this work.

\bibliographystyle{plainnat}
\bibliography{references}

\end{document}
